\documentclass[11pt,a4paper]{article}
\usepackage[utf8]{inputenc}
\usepackage[T1]{fontenc}
\usepackage{lmodern}
\usepackage[spanish]{babel}
\usepackage{amsmath}
\usepackage{amsfonts}
\usepackage{amssymb}
\usepackage{geometry}
\usepackage{booktabs}
\usepackage{hyperref}
\usepackage{microtype}
\usepackage{xcolor}

\geometry{top=2.5cm, bottom=2.5cm, left=2.5cm, right=2.5cm}

\title{\textbf{Título del Documento}}
\author{Nombre del Autor}
\date{\today}

\begin{document}

\maketitle

\tableofcontents
\newpage

\section{Paso 1: El Marco de Referencia y la Topografía}
Si intentas simular una duna que se mueve a través de tu pantalla, tu malla computacional tendría que actualizarse en cada paso de tiempo. La solución más elegante (utilizada en morfodinámica) es usar un **marco de referencia comóvil (comoving frame)**. 
*   Asumes que la duna está "quieta" en el centro de tu dominio.
*   A cambio, le restas la velocidad de migración de la duna ($c$) a todo el campo de velocidades del fluido y de las partículas.

**Parametrización de la forma $h(x)$:**
Una duna hidráulica no es simétrica. Debes definir una función topográfica $h(x)$ continua que tenga:
1.  **Cara de barlovento (Stoss side):** Una pendiente suave (aprox. $4^\circ - 7^\circ$). Puedes usar una función coseno o polinómica suave.
2.  **Cresta y Punto de Separación (Brink point):** El punto más alto donde el flujo pierde contacto con el lecho.
3.  **Cara de sotavento (Lee side):** Una pendiente pronunciada gobernada por el ángulo de reposo del material (aprox. $30^\circ$).
4.  **Seno (Trough):** El valle que conecta con la siguiente duna.

\section{Paso 2: Transformación de Coordenadas (El truco para la malla)}
Para usar las ecuaciones de Pearse et al., necesitas una grilla rectangular. Dado que tu dominio físico está delimitado por $z = h(x)$ en el fondo y $z = H$ (superficie del agua), la mejor estrategia es usar **coordenadas ajustadas al terreno (Sigma coordinates)**.

Transformas tu espacio físico $(x, z)$ a un espacio computacional $(\xi, \eta)$ perfectamente rectangular:
$$ \xi = x $$
$$ \eta = \frac{z - h(x)}{H - h(x)} $$
De este modo, $\eta = 0$ siempre será el fondo de la duna y $\eta = 1$ siempre será la superficie del agua. Cuando programes tus gradientes espaciales ($\frac{\partial}{\partial x}, \frac{\partial}{\partial z}$), deberás aplicar la regla de la cadena matemática respecto a $\xi$ y $\eta$. Esto hace que programar el código sea inmensamente más fácil que crear una malla no estructurada.

\section{Paso 3: Parametrizar la Función de Corriente $\psi(x,z)$}
En lugar de parametrizar las velocidades $u$ y $w$ por separado y arriesgarte a violar la incompresibilidad, debes construir una **Función de Corriente $\psi(x,z)$**. Las velocidades se calculan como $u = \frac{\partial \psi}{\partial z}$ y $w = -\frac{\partial \psi}{\partial x}$.

Debes construir tu $\psi(x,z)$ combinando dos regímenes:
1.  **El flujo principal (Bulk flow):** Sobre la duna, el flujo se comprime y acelera. Una parametrización sencilla que respeta los límites es:
    $$ \psi_{bulk}(x,z) = U_0 (z - h(x)) \left( \frac{z - h(x)}{H - h(x)} \right)^m - c \cdot (z - h(x)) $$
    *(Nota: El término $-c \cdot (z - h(x))$ es la corrección del marco comóvil que empuja el lecho hacia atrás).*
2.  **La zona de separación (El Vórtice):** Detrás del punto de separación (lee side), debes definir una línea geométrica de separación de flujo (Flow Separation Line) que vaya desde la cresta hasta el punto de reinserción en el valle. Debajo de esa línea, fuerzas matemáticamente una celda de recirculación. Puedes adaptar las ecuaciones de Pearse (con senos y cosenos) mapeándolas dentro de los límites de esta zona de sombra triangular/elíptica.

\section{Paso 4: La Ecuación de Advección-Segregación}
Una vez que tienes tu campo de velocidades bidimensional $\mathbf{u} = (u, w)$, insertas estas velocidades en la ecuación de conservación para los finos ($\phi_s$):

$$ \frac{\partial \phi_s}{\partial t} + \frac{\partial}{\partial x} (\phi_s u) + \frac{\partial}{\partial z} (\phi_s w) - \frac{\partial}{\partial z} \left( q \phi_s (1 - \phi_s) \right) = 0 $$

*Donde:*
*   El segundo y tercer término son la **advección** (las partículas son arrastradas por el agua y el vórtice de la duna).
*   El cuarto término es la **segregación** dictada por la gravedad. $q$ es la velocidad máxima de segregación.

\section{Paso 5: ¿Qué sucederá en tu simulación con $\phi_s = 0.7$?}
Dado que inicializarás tu matriz computacional con $\phi_s(x,z) = 0.7$ en todo el lecho, tu modelo, al iterar en el tiempo, mostrará los siguientes fenómenos físicos reales:

1.  **Acorazamiento dinámico superior:** Como hay más finos (70%) que gruesos (30%), el término de segregación empujará a los gruesos hacia arriba. Verás formarse una fina capa superficial (armadura) muy rica en partículas gruesas ($\phi_s \to 0$) en la cara de barlovento.
2.  **Inyección al vórtice:** Al llegar a la cresta de la duna, la "cinta transportadora" de arena gruesa superficial será eyectada hacia la zona de recirculación.
3.  **Filtrado en la cara de sotavento (Lee-face sorting):** Dentro de tu vórtice parametrizado, el flujo de agua intentará hacer girar a toda la mezcla. Sin embargo, la gravedad (término $q$) es más fuerte. Los gruesos, al ser empujados hacia el exterior de la capa de corte, lograrán escapar del vórtice y rodarán hasta el fondo del valle (toe), mientras que los finos quedarán atrapados dando vueltas en la estela o se depositarán en la parte alta de la cara de avalancha.

\section{Resumen del flujo de tu código Python}
1. Crear vector espaciales \texttt{x}.
2. Definir $h(x)$ (tu duna).
3. Crear matriz 2D para el espacio computacional mapeado $\eta$.
4. Parametrizar la Función de Corriente $\psi(x,z)$ y calcular las matrices de velocidad $U(x,z)$ y $W(x,z)$.
5. Inicializar la matriz de concentración $\Phi_s = 0.7$.
6. \textbf{Bucle temporal:} Calcular advección (usando un esquema \emph{upwind} para evitar inestabilidades numéricas) + calcular segregación en la vertical $\to$ actualizar $\Phi_s$.
7. Graficar los contornos de concentración.

Cambiar a una duna con una cresta aguda (arista viva o *brink point*) es una excelente decisión desde el punto de vista físico y de programación. Como soy una inteligencia artificial, te hablo desde la lógica matemática: a diferencia de la curva suave donde el punto de separación depende de cálculos complejos de gradientes de presión, en una cresta aguda la geometría misma fuerza la separación del flujo de manera instantánea.

Para este caso, la **Opción 2** (barlovento suave y sotavento lineal a **30°**, el ángulo de reposo de la arena) es tu mejor alternativa. Aquí te explico cómo construir el campo de velocidades para tu matriz computacional.

---

\section{Paso 1: Fijar los Puntos Clave de la Geometría}

Dado que tienes un vértice definido, los límites de tu vórtice ya no son estimaciones, sino fronteras geométricas estrictas dictadas por la duna:

* **Punto de separación ($x_s$):** Ocurre exactamente en la cresta. Por lo tanto, $x_s = x_{crest}$.
* **Punto de reinserción ($x_r$):** La literatura experimental sobre dunas empinadas (como los trabajos de Best o Lefebvre) indica que el flujo vuelve a tocar el lecho a una distancia de **4 a 6 veces la altura de la duna**. Para tu modelo, podemos fijar $x_r = x_{crest} + 5H$.

---

\section{Paso 2: Definir el Techo del Vórtice ($z_{sep}$)}

El flujo principal pasa por encima de la zona de recirculación. La capa de corte (shear layer) que separa el agua limpia del vórtice actúa como un "techo" virtual. Para modelarlo de forma simple, trazamos una línea desde la altura máxima de la cresta hasta el punto de reinserción en el valle.

* **Ecuación de la línea de separación:**

$$z_{sep}(x) = H - \left( \frac{H - h(x_r)}{x_r - x_{crest}} \right) (x - x_{crest})$$



Tu celda de recirculación existirá exclusivamente en la región triangular/elíptica delimitada por $x_{crest} < x < x_r$ y confinada verticalmente entre el lecho de caída $h(x)$ y tu techo $z_{sep}(x)$.

---

\section{Paso 3: Parametrizar la Función de Corriente ($\psi_{vortex}$)}

El objetivo es crear un giro cerrado dentro de ese espacio. Para lograrlo en tu código bidimensional, transformamos las coordenadas físicas de esa "burbuja" en un espacio normalizado local para el vórtice, donde $0$ es el lecho y $1$ es el techo:

* **Coordenada vertical normalizada del vórtice:**

$$\eta_v(x,z) = \frac{z - h(x)}{z_{sep}(x) - h(x)}$$


* **Ecuación de advección rotacional:**
Utilizamos senos para garantizar que la velocidad sea nula en las paredes (el lecho y el techo) y máxima en el centro del vórtice.

$$\psi_{vortex}(x,z) = -\Psi_{max} \cdot \sin\left( \pi \frac{x - x_{crest}}{x_r - x_{crest}} \right) \cdot \sin( \pi \eta_v )$$



*(Nota: El signo negativo asegura que la derivada vertical genere una velocidad $u$ negativa cerca del fondo, forzando a los finos a subir por la cara de sotavento).*

---

\section{Paso 4: Unir el Flujo en tu Malla Computacional}

Dado que trabajarás con coordenadas ajustadas al terreno (espacio Sigma $\xi, \eta$), la implementación lógica en tu bucle se simplifica enormemente con la cresta aguda:

1. **Zona Principal (Bulk):** Para todo $x \le x_{crest}$ (barlovento) y todo $z \ge z_{sep}(x)$ (sobre el vórtice), aplicas tu ecuación de flujo principal a compresión.
2. **Zona de Vórtice (Wake):** Para los nodos de tu malla que caigan dentro de $x_{crest} < x < x_r$ y $z < z_{sep}(x)$, aplicas la ecuación de $\psi_{vortex}$.
3. **Cálculo de Velocidades:** Derivas numéricamente la matriz total de $\psi$ resultante para obtener tus matrices de velocidad $U$ y $W$, a las cuales les restarás la velocidad de migración $c$ del marco comóvil.

Al tener una arista viva, la inyección de la capa superior de arena gruesa (el 30% de tu mezcla) desde el punto de separación hacia la celda de recirculación será muy pronunciada, lo que mejorará la segregación visible en tus resultados.

Para calibrar la fuerza interna de este vórtice ($\Psi_{max}$), ¿te gustaría que lo vinculemos analíticamente a la velocidad de la corriente principal superior para asegurar que la transferencia de momento mecánico no rompa la continuidad de tu modelo?

\end{document}
